\documentclass[aspectratio=169,11pt]{beamer}
\usepackage{booktabs}
\usepackage{pgfplots}
\usepackage{pgfplotstable}
\usepackage{amsmath}
\usepackage{array}
\usepackage[backend=biber,style=authoryear]{biblatex}
\addbibresource{references.bib}

\pgfplotsset{compat=1.18}
% ============================================================
% Synthetic Library Space Utilisation Dataset
% ============================================================
%
% This dataset is fully synthetic and was constructed specifically
% for this benchmark asset. It is NOT presented as real university data.
%
% Observation unit:
% One study zone x one time period x one day.
%
% Observation period:
% One hypothetical week, Monday to Saturday.
%
% Study zones and capacities:
% Quiet Study Zone   = 40 seats
% General Study Zone = 60 seats
% Group Study Zone   = 30 seats
%
% Time periods:
% P1 = 8:00-10:00 AM
% P2 = 10:00 AM-12:00 PM
% P3 = 2:00-4:00 PM
% P4 = 4:00-6:00 PM
%
% Occupancy rate is derived as:
% OccupiedSeats / Capacity * 100
% ============================================================

\pgfplotstableread[col sep=comma]{
Day,Period,Zone,Capacity,OccupiedSeats,AvgStudyDuration,StudentEntries,SatisfactionScore
Monday,P1,Quiet,40,28,85,24,3.3
Monday,P1,General,60,38,93,34,3.5
Monday,P1,Group,30,16,79,9,3.3
Monday,P2,Quiet,40,31,92,17,3.1
Monday,P2,General,60,46,51,30,3.0
Monday,P2,Group,30,19,74,14,3.3
Monday,P3,Quiet,40,28,78,18,3.2
Monday,P3,General,60,44,102,36,3.4
Monday,P3,Group,30,17,102,9,3.8
Monday,P4,Quiet,40,37,80,19,3.0
Monday,P4,General,60,50,101,33,3.3
Monday,P4,Group,30,21,53,10,2.3
Tuesday,P1,Quiet,40,29,89,27,2.9
Tuesday,P1,General,60,41,96,30,3.3
Tuesday,P1,Group,30,16,80,15,3.4
Tuesday,P2,Quiet,40,32,69,27,2.9
Tuesday,P2,General,60,49,61,34,3.0
Tuesday,P2,Group,30,21,106,14,3.8
Tuesday,P3,Quiet,40,31,62,21,2.5
Tuesday,P3,General,60,42,81,27,3.2
Tuesday,P3,Group,30,18,47,9,2.2
Tuesday,P4,Quiet,40,36,52,26,3.0
Tuesday,P4,General,60,55,99,45,3.6
Tuesday,P4,Group,30,22,49,12,2.9
Wednesday,P1,Quiet,40,29,64,23,2.5
Wednesday,P1,General,60,41,58,32,3.1
Wednesday,P1,Group,30,17,80,16,3.0
Wednesday,P2,Quiet,40,34,57,17,2.6
Wednesday,P2,General,60,48,47,36,2.6
Wednesday,P2,Group,30,19,103,11,3.6
Wednesday,P3,Quiet,40,31,76,21,3.1
Wednesday,P3,General,60,46,49,28,3.0
Wednesday,P3,Group,30,17,53,17,3.1
Wednesday,P4,Quiet,40,37,91,24,3.6
Wednesday,P4,General,60,56,72,56,3.0
Wednesday,P4,Group,30,23,99,21,3.1
Thursday,P1,Quiet,40,33,79,22,3.4
Thursday,P1,General,60,36,58,36,3.1
Thursday,P1,Group,30,18,54,8,3.1
Thursday,P2,Quiet,40,36,83,27,3.3
Thursday,P2,General,60,47,76,28,2.9
Thursday,P2,Group,30,19,75,14,2.9
Thursday,P3,Quiet,40,30,74,20,3.4
Thursday,P3,General,60,41,74,27,2.8
Thursday,P3,Group,30,17,77,9,3.4
Thursday,P4,Quiet,40,38,60,26,3.1
Thursday,P4,General,60,52,75,27,3.5
Thursday,P4,Group,30,24,73,17,3.5
Friday,P1,Quiet,40,29,72,17,3.4
Friday,P1,General,60,37,67,31,3.3
Friday,P1,Group,30,15,59,15,3.0
Friday,P2,Quiet,40,33,67,16,2.7
Friday,P2,General,60,44,46,45,2.5
Friday,P2,Group,30,20,101,13,3.5
Friday,P3,Quiet,40,28,99,26,3.4
Friday,P3,General,60,37,100,28,3.1
Friday,P3,Group,30,18,61,8,2.6
Friday,P4,Quiet,40,36,106,17,3.7
Friday,P4,General,60,50,107,47,3.4
Friday,P4,Group,30,22,91,16,3.3
Saturday,P1,Quiet,40,27,72,22,3.1
Saturday,P1,General,60,30,48,14,2.7
Saturday,P1,Group,30,12,107,8,4.0
Saturday,P2,Quiet,40,32,92,30,3.7
Saturday,P2,General,60,39,56,18,2.6
Saturday,P2,Group,30,18,75,8,3.0
Saturday,P3,Quiet,40,27,89,28,3.1
Saturday,P3,General,60,40,106,34,3.3
Saturday,P3,Group,30,15,76,16,3.3
Saturday,P4,Quiet,40,33,72,23,3.1
Saturday,P4,General,60,50,108,29,3.3
Saturday,P4,Group,30,19,45,15,2.9
}\librarydata
\usetheme{Madrid}

\title{Analysing University Library Space Utilisation and Student Study Patterns}
\author{Oluwadurotimi Alade}
\institute{University of Ibadan}
\date{September 2026}

\begin{document}

\begin{frame}
    \titlepage
\end{frame}

\begin{frame}{Presentation Overview}
    \tableofcontents
\end{frame}

\section{Introduction}

\begin{frame}{Background}
    \begin{itemize}
        \item University libraries provide spaces for individual and collaborative study.
        \item The way students use these spaces can vary across zones, time periods, and days.
        \item Understanding these patterns can support better use of available study spaces.
    \end{itemize}
\end{frame}

\begin{frame}{Problem Statement}
    Library study spaces have limited seating capacity, but demand may vary
    across different zones, time periods, and days. A structured analysis of
    occupancy, study duration, student entries, and satisfaction can help
    identify patterns in space utilisation.
\end{frame}

\begin{frame}{Aim and Objectives}
    \textbf{Aim}

    To analyse university library space utilisation and student study patterns
    using synthetic observational data.

    \vspace{0.4cm}

    \textbf{Objectives}
    \begin{enumerate}
        \item Measure occupancy rates across library study zones.
        \item Compare utilisation across time periods and days.
        \item Examine study duration, student entries, and satisfaction.
        \item Explore relationships among selected study-pattern variables.
        \item Develop practical recommendations from the observed patterns.
    \end{enumerate}
\end{frame}

\section{Research Questions}

\begin{frame}{Research Questions}
    \begin{enumerate}
        \item How does library space occupancy vary across the three study zones?
        \item How does occupancy vary across different time periods?
        \item How does occupancy vary across days of the week?
        \item What patterns exist between student entries, study duration, and occupancy?
        \item What relationship exists between occupancy and student satisfaction?
    \end{enumerate}
\end{frame}

\section{Literature Review}

\begin{frame}{Literature Review}
	\begin{itemize}
		\item Physical library space can influence how students choose and use study areas \parencite{cha2015physical}.
		
		\item Students use academic library spaces for different individual and collaborative study activities \parencite{bedwell2013seeing}.
		
		\item Students' needs and preferences can vary across different library spaces and study environments \parencite{may2015using}.
	\end{itemize}
\end{frame}

\begin{frame}{Conceptual Background}
	\begin{itemize}
		\item Library space utilisation reflects how available study spaces are occupied and used by students.
		
		\item Students may select different study environments depending on their study activities and preferences \parencite{cha2015physical}.
		
		\item Academic library spaces can support both individual and collaborative study activities \parencite{bedwell2013seeing}.
		
		\item The design and characteristics of study spaces can influence students' experience of the library environment \parencite{defrain2020interiors,hegde2018work}.
	\end{itemize}
\end{frame}

\begin{frame}{Dataset Overview}
	\begin{table}
		\centering
		\begin{tabular}{lr}
			\toprule
			\textbf{Dataset Feature} & \textbf{Value} \\
			\midrule
			Observation days & 6 \\
			Observation periods per day & 4 \\
			Study zones & 3 \\
			Total observations & 72 \\
			\bottomrule
		\end{tabular}
	\end{table}
	
	\vspace{0.4cm}
	
	\begin{itemize}
		\item The dataset covers a hypothetical Monday-to-Saturday observation week.
		\item Each observation represents one study zone, one period, and one day.
		\item The dataset is fully synthetic and embedded in the project.
	\end{itemize}
\end{frame}

\begin{frame}{Variables and Data Structure}
	\small
	\begin{table}
		\centering
		\begin{tabular}{p{3.2cm}p{6.8cm}}
			\toprule
			\textbf{Variable} & \textbf{Description} \\
			\midrule
			Day & Day of observation \\
			Period & Two-hour observation period \\
			Zone & Library study zone \\
			Capacity & Number of available seats \\
			Occupied Seats & Number of occupied seats during observation \\
			Avg. Study Duration & Average study duration in minutes \\
			Student Entries & Number of student entries recorded \\
			Satisfaction Score & Student satisfaction rating on a 1--5 scale \\
			\bottomrule
		\end{tabular}
	\end{table}
	
	\vspace{0.2cm}
	
	\begin{itemize}
		\item Occupancy rate is derived from Occupied Seats and Capacity.
		\item The raw observations are embedded in \texttt{data.tex}.
	\end{itemize}
\end{frame}

\section{Data and Methodology}

\begin{frame}{Study Design}
    \begin{itemize}
        \item Data type: fully synthetic benchmark data.
        \item Observation period: one hypothetical week, Monday to Saturday.
        \item Study zones: Quiet, General, and Group.
        \item Four two-hour observation periods were used each day.
        \item Total observations: 72.
    \end{itemize}
\end{frame}

\begin{frame}{Variables}
    \begin{table}
        \centering
        \small
        \begin{tabular}{ll}
            \toprule
            \textbf{Variable} & \textbf{Description} \\
            \midrule
            Day & Day of observation \\
            Period & Two-hour observation period \\
            Zone & Library study zone \\
            Capacity & Available seats in the zone \\
            Occupied Seats & Seats occupied during observation \\
            Study Duration & Average study duration in minutes \\
            Student Entries & Student entries recorded during the period \\
            Satisfaction & Average satisfaction score from 1 to 5 \\
            \bottomrule
        \end{tabular}
    \end{table}
\end{frame}

\begin{frame}{Observation Periods}
    \begin{table}
        \centering
        \small
        \begin{tabular}{cl}
            \toprule
            \textbf{Period} & \textbf{Time} \\
            \midrule
            P1 & 8:00--10:00 AM \\
            P2 & 10:00 AM--12:00 PM \\
            P3 & 2:00--4:00 PM \\
            P4 & 4:00--6:00 PM \\
            \bottomrule
        \end{tabular}
    \end{table}

    \vspace{0.3cm}

    Each observation covers one study zone during one two-hour period
    on one day.
\end{frame}

\section{Preliminary Results}

\begin{frame}{Preliminary Occupancy Result}
    The current validated dataset gives an overall mean occupancy rate of:

    \[
    \boxed{71.59\%}
    \]
\end{frame}

\section{Methodology}

\begin{frame}{Research Design and Methodology}
	\begin{itemize}
		\item The study uses a descriptive and comparative observational design.
		
		\item The observation period covers a hypothetical Monday-to-Saturday week.
		
		\item Observations were recorded across three study zones and four two-hour periods each day.
		
		\item Each observation represents one study zone, one period, and one day.
		
		\item The analysis focuses on occupancy, study duration, student entries, and satisfaction.
		
		\item All observations are synthetic and embedded in the project for benchmark reproducibility.
	\end{itemize}
\end{frame}

\begin{frame}{Statistical Analysis Approach}
	\begin{itemize}
		\item Occupancy rate is calculated as:
		\[
		\text{Occupancy Rate} =
		\frac{\text{Occupied Seats}}{\text{Capacity}}\times100
		\]
		
		\item Descriptive statistics are used to summarise occupancy, study duration, student entries, and satisfaction.
		
		\item Mean occupancy is compared across study zones, observation periods, and days.
		
		\item Pearson correlation is used to examine linear associations between the quantitative variables.
		
		\item Results are interpreted as descriptive associations rather than causal effects.
	\end{itemize}
\end{frame}

\begin{frame}{Descriptive Analysis}
	\small
	\begin{table}
		\centering
		\begin{tabular}{lr}
			\toprule
			\textbf{Measure} & \textbf{Overall Mean} \\
			\midrule
			Occupancy Rate & 71.59\% \\
			Average Study Duration & 76.53 minutes \\
			Student Entries & 22.60 \\
			Satisfaction Score & 3.14 / 5 \\
			\bottomrule
		\end{tabular}
	\end{table}
	
	\vspace{0.3cm}
	
	\begin{itemize}
		\item Mean occupancy across the 72 observations was 71.59\%.
		\item Average study duration was 76.53 minutes.
		\item The mean number of student entries was 22.60.
		\item Mean satisfaction was 3.14 on the 1--5 scale.
	\end{itemize}
\end{frame}

\begin{frame}{Comparative and Correlation Analysis}
	\scriptsize
	
	\begin{columns}[T]
		
		\begin{column}{0.48\textwidth}
			\textbf{Mean Occupancy by Zone}
			
			\begin{tabular}{lr}
				\toprule
				\textbf{Zone} & \textbf{Occupancy} \\
				\midrule
				Quiet & 79.69\% \\
				General & 73.54\% \\
				Group & 61.53\% \\
				\bottomrule
			\end{tabular}
			
			\vspace{0.25cm}
			
			\textbf{Mean Occupancy by Period}
			
			\begin{tabular}{lr}
				\toprule
				\textbf{Period} & \textbf{Occupancy} \\
				\midrule
				P1 & 62.36\% \\
				P2 & 74.26\% \\
				P3 & 66.34\% \\
				P4 & 83.38\% \\
				\bottomrule
			\end{tabular}
		\end{column}
		
		\begin{column}{0.48\textwidth}
			\textbf{Pearson Correlations}
			
			\begin{tabular}{lr}
				\toprule
				\textbf{Variable Pair} & \textbf{r} \\
				\midrule
				Occupancy--Duration & 0.063 \\
				Occupancy--Entries & 0.508 \\
				Occupancy--Satisfaction & -0.013 \\
				Duration--Entries & 0.104 \\
				Duration--Satisfaction & 0.726 \\
				Entries--Satisfaction & 0.018 \\
				\bottomrule
			\end{tabular}
			
			\vspace{0.25cm}
			
			\begin{itemize}
				\item P4 had the highest mean occupancy.
				\item Occupancy and student entries showed a moderate positive association.
				\item Study duration and satisfaction showed a strong positive association.
			\end{itemize}
		\end{column}
		
	\end{columns}
\end{frame}

\section{Results and Discussion}

\begin{frame}{Results and Discussion}
	\small
	\begin{itemize}
		\item The Quiet Zone recorded the highest mean occupancy at 79.69\%, followed by the General Zone at 73.54\% and the Group Zone at 61.53\%.
		
		\item P4 recorded the highest mean occupancy at 83.38\%, indicating greater space use during the final observation period.
		
		\item Saturday recorded the lowest mean daily occupancy at 64.65\%, while Wednesday recorded the highest at 74.93\%.
		
		\item Occupancy and student entries showed a moderate positive association ($r=0.508$).
		
		\item Study duration and satisfaction showed a strong positive association ($r=0.726$).
		
		\item These findings describe patterns and associations within the synthetic benchmark dataset and do not establish causal relationships.
	\end{itemize}
\end{frame}

\begin{frame}{Recommendations and Limitations}
	\small
	
	\textbf{Recommendations}
	\begin{itemize}
		\item Review space allocation in the Quiet Zone because it recorded the highest mean occupancy.
		\item Pay particular attention to demand during P4, which recorded the highest mean occupancy.
		\item Consider study duration when evaluating student experience, given its positive association with satisfaction.
		\item Continue monitoring occupancy and student use across individual zones and periods.
	\end{itemize}
	
	\vspace{0.15cm}
	
	\textbf{Limitations}
	\begin{itemize}
		\item The dataset is fully synthetic and represents one hypothetical observation week.
		\item The analysis covers three study zones and four observation periods per day.
		\item The observational design identifies patterns and associations but does not establish causation.
	\end{itemize}
	
\end{frame}

\begin{frame}{Conclusion}
	\small
	\begin{itemize}
		\item The analysis shows clear differences in library space utilisation across study zones and observation periods.
		
		\item The Quiet Zone recorded the highest mean occupancy, while the Group Zone recorded the lowest.
		
		\item P4 was the busiest observation period, while Saturday had the lowest mean daily occupancy.
		
		\item Occupancy was moderately associated with student entries, while study duration showed a strong positive association with satisfaction.
		
		\item The findings provide a basis for targeted monitoring and planning of library study spaces.
	\end{itemize}
	
	\vspace{0.35cm}
	
	\centering
	\textbf{Thank you}
	
\end{fram}

\begin{frame}
    \centering
    \Huge Questions?
\end{frame}

\end{document}
